\documentclass[a4paper]{article}

\usepackage[T1]{fontenc}
\usepackage[utf8]{inputenc}
\usepackage{mlmodern}

%\usepackage{ngerman}	% Sprachanpassung Deutsch

\usepackage{graphicx}
\usepackage{geometry}
\geometry{a4paper, top=15mm}

\usepackage{subcaption}
\usepackage[shortlabels]{enumitem}
\usepackage{amssymb}
\usepackage{amsthm}
\usepackage{amsmath}
\usepackage{mathtools}
\usepackage{braket}
\usepackage{bbm}
\usepackage{graphicx}
\usepackage{float}
\usepackage{yhmath}
\usepackage{tikz}
\usepackage{scratch}
\usepackage{algorithm}
\usepackage{algpseudocode}
\usetikzlibrary{patterns,decorations.pathmorphing,positioning}
\usetikzlibrary{calc,decorations.markings, shapes.geometric}

\usepackage[backend=biber, sorting=none]{biblatex}
\addbibresource{../cite.bib}

\usepackage[framemethod=TikZ]{mdframed}

\tikzstyle{titlered} =
    [draw=black, thick, fill=white,%
        text=black, rectangle,
        right, minimum height=.7cm]


\usepackage[colorlinks=true,naturalnames=true,plainpages=false,pdfpagelabels=true]{hyperref}

\usepackage[OT2,T1]{fontenc}

\title{
    Exposé\\
}
\author{Milutin Popović}
\date{}


\begin{document}

\maketitle
\thispagestyle{empty}
\vspace*{\fill}
The thesis deals with the analysis and the development of a solution to the
arbitrage problem in decentralized finance. In this context, arbitrage is a
subform of strategies used to extract Maximal Extractable Value (MEV) on
public blockchains, specifically with smart contract functionality.

MEV sometimes also called Miner Extractable Value or Block Proposer
Extractable Value refers to the maximal value that can be extracted from
block production, without considering block rewards or gas fees. This is done
by including or excluding transactions in the public transaction pool and
combining them with custom transactions specifically aimed to exploit the
market inefficiency. Due to the short time window for block production, e.g.
on the Ethereum chain blocks are included every 12 seconds, extremely
efficient methods are needed for extracting MEV. While some forms of MEV,
essentially exploit the transaction of an end user and give them an overall
worse-than-expected outcome, others like the arbitrage do not and are crucial
to the system by providing market stability.

In this scope, the thesis will first formalize the building blocks of
decentralized finance, which are the automated market makers, and together with
this frame the nonlinear optimization problem resembling the arbitrage opportunity. The
nonlinear optimization problem will be mathematically analyzed and based on the
analysis an algorithm will be proposed for efficient computation of the
solution. In the pre-thesis setting while brainstorming about the solutions to
the problem an interesting question emerged: How to spot an arbitrage
opportunity?

To answer this question in the time constraints given by the block building
process and the rate of incoming transactions in the transaction pool, a
reinforcement learning method is utilized. The agent will learn the policy
and value functions of the Markov Decision Process by using deep graph neural
networks with attention. The environment will resemble a graph with nodes
being on-chain tokens and edges being on-chain pools. This formulation allows
the question of finding the arbitrage opportunity to be framed as a
combinatorial optimization problem. The agent will learn to navigate this
graph to find a profitable arbitrage walk. The walk is propagated to
construct the optimization problem and solve it.
\vspace*{\fill}

\end{document}
